% Projections used to look at the latent space. Subsection of "Clustering".

\subsection{Visualizing the latent space}  % kernel PCA, t-SNE, UMAP
\label{sec:clust_projections}            % \secref{sec:clust_projections}

The latent codes are represented in a 256-dimensional space and therefore cannot be directly
visualized. We project them to two and three dimensions using three complementary dimensionality-
reduction methods and colour the resulting embeddings according to cluster membership or
morphological descriptors.
% The projections are used only for viewing: all clusters are computed in the full latent space and no projection preserves all
% distances, so a separation that appears in a projection, or the lack of one, has to be read with
% care. 
The three methods distort the data in different ways, which is why we compare them.

\paragraph{Kernel PCA}                   % linear PCA in the RBF feature space
Principal component analysis (PCA) finds the orthogonal directions of largest variance; kernel PCA
\cite{scholkopf1998kpca} does the same in the feature space of a kernel. With the RBF kernel of
Eq.~\eqref{eq:rbf}
\begin{equation}
  k(\mathbf{x}_i,\mathbf{x}_j) = \exp\!\left(-\gamma \|\mathbf{x}_i-\mathbf{x}_j\|^2\right)
  \label{eq:rbf}
\end{equation}
and the same median-rule $\gamma$ as kernel K-means, it shows the space in
which kernel K-means operates. The kernel matrix $K$ is first centred,
$\tilde{K} = H K H$ with $H = I - \frac{1}{N}\mathbf{1}\mathbf{1}^{\top}$, and decomposed as
$\tilde{K}\mathbf{v}_m = \lambda_m \mathbf{v}_m$. The coordinate of tree $i$ on the $m$-th
component is $\sqrt{\lambda_m}\, v_{m,i}$, and the share of kernel variance kept by a view with
$d$ components is $\sum_{m \le d} \lambda_m / \sum_m \lambda_m$. 
% Kernel PCA is deterministic and
% keeps the global layout of the data, but it only shows the few directions of largest variance.
% % \TODO{Report the kernel variance shown in 2D and in 3D (printed by the notebook).}

\paragraph{t-SNE}                        % neighbourhood probabilities, KL divergence
t-distributed stochastic neighbour embedding (t-SNE) \cite{vandermaaten2008tsne} preserves
neighbourhoods rather than distances. For every tree $i$, it turns distances into probabilities
of picking tree $j$ as a neighbour,
\begin{equation}
  p_{j|i} = \frac{\exp\!\big(-\|\mathbf{x}_i - \mathbf{x}_j\|^2 / 2\sigma_i^2\big)}
                 {\sum_{k \neq i} \exp\!\big(-\|\mathbf{x}_i - \mathbf{x}_k\|^2 / 2\sigma_i^2\big)} ,
  \label{eq:tsne_p}                      % neighbour probabilities in the latent space
\end{equation}
where the width $\sigma_i$ is chosen so that the distribution has a given \emph{perplexity},
roughly the effective number of neighbours. In the projection, similarities $q_{ij}$ follow a
heavy-tailed Student-$t$ distribution, $q_{ij} \propto (1 + \|\mathbf{y}_i - \mathbf{y}_j\|^2)^{-1}$,
and the points $\mathbf{y}_i$ are moved by gradient descent to minimize the Kullback--Leibler
divergence between the symmetrized $p_{ij}$ and $q_{ij}$. The heavy tail lets dissimilar trees
drift apart, which gives well-separated groups, but the sizes of the groups and the distances
between them are not meaningful. We ran t-SNE with PCA initialization and 1000 iterations, for
perplexities from 5 to 50 in steps of 5.

\paragraph{UMAP}
Uniform Manifold Approximation and Projection (UMAP) \cite{mcinnes2018umap} is a nonlinear
dimensionality reduction method that seeks to preserve the neighbourhood structure of the
high-dimensional data in a lower-dimensional representation. Given the latent codes
$\mathbf{x}_1,\ldots,\mathbf{x}_N$, UMAP first identifies the nearest neighbours of each
observation using the Euclidean distance. These neighbourhood relationships are represented by
a weighted graph.

For each tree $i$, UMAP determines a local connectivity distance $\rho_i$, corresponding to the
distance to the closest neighbour, and a local scale parameter $\sigma_i$. The membership strength
of the connection between tree $i$ and one of its neighbours $j$ is defined as
\begin{equation}
    w_{ij} =
    \begin{cases}
        1, & d(\mathbf{x}_i,\mathbf{x}_j) \leq \rho_i,\\[4pt]
        \exp\!\left(
        -\dfrac{d(\mathbf{x}_i,\mathbf{x}_j)-\rho_i}{\sigma_i}
        \right),
        & d(\mathbf{x}_i,\mathbf{x}_j) > \rho_i,
    \end{cases}
    \label{eq:umap_weight}
\end{equation}
where $d(\mathbf{x}_i,\mathbf{x}_j)$ denotes the Euclidean distance between two latent codes.
Thus, nearby observations receive a high connection strength, while more distant neighbours have
smaller weights. The local scales are chosen so that the graph adapts to the density of the data
around each observation.

UMAP then seeks a low-dimensional representation
$\mathbf{y}_1,\ldots,\mathbf{y}_N$ whose neighbourhood structure reproduces that of the original
graph. In the low-dimensional space, the strength of an association between two observations is
modelled by
\begin{equation}
    \tilde{w}_{ij}
    =
    \frac{1}
    {1+a\|\mathbf{y}_i-\mathbf{y}_j\|_2^{2b}},
    \label{eq:umap_lowdim}
\end{equation}
where $a$ and $b$ determine the shape of the low-dimensional decay function. The embedding is
obtained by minimizing the cross-entropy between the high- and low-dimensional fuzzy graphs:
\begin{equation}
    \mathcal{L}_{\mathrm{UMAP}}
    =
    -\sum_{i,j}
    \left[
        w_{ij}\log(\tilde{w}_{ij})
        +
        (1-w_{ij})\log(1-\tilde{w}_{ij})
    \right].
    \label{eq:umap_loss}
\end{equation}
Minimizing this objective attracts observations that are strongly connected in the original
latent space and repels observations that have weak or no connections. The resulting embedding
therefore preserves neighbourhood relationships rather than the exact pairwise distances of the
original space.

The scale of the neighbourhood structure is controlled by the number of neighbours considered for
each observation. Small values emphasize local structure, whereas larger values allow broader
relationships to influence the embedding. A minimum-distance parameter further controls how
closely neighbouring observations can be placed in the low-dimensional space.

In our experiments, UMAP was applied directly to the 256-dimensional latent representations using
Euclidean distance and the default minimum-distance value. The neighbourhood size was varied
from 5 to 55 in increments of 5, and both two- and three-dimensional embeddings were generated.

% \paragraph{Parameter sweeps}             % trust what is stable across settings
% Because t-SNE and UMAP depend on their parameters and on the random initialization (fixed with a
% seed), we computed them over the whole range of perplexities and neighbour counts, in 2D and 3D.
% A grouping that appears for every setting reflects the data; one that appears for a single
% setting is more likely an artefact of the projection.
